\documentclass[12pt,a4paper,bibliography=totocnumbered,listof=totocnumbered]{scrartcl}
\usepackage[ngerman]{babel}
\usepackage[utf8]{inputenc}
\usepackage{amsmath}
\usepackage{amsfonts}
\usepackage{amssymb}
\usepackage{graphicx}
\usepackage{fancyhdr}
\usepackage{tabularx}
\usepackage{geometry}
\usepackage{setspace}
\usepackage[right]{eurosym}
\usepackage[printonlyused]{acronym}
\usepackage{subfig}
\usepackage{floatflt}
\usepackage[usenames,dvipsnames]{color}
\usepackage{colortbl}
\usepackage{paralist}
\usepackage{array}
\usepackage{titlesec}
\usepackage{parskip}
\usepackage[right]{eurosym}
\usepackage{picins}
\usepackage[subfigure,titles]{tocloft}
\usepackage[pdfpagelabels=true]{hyperref}

\usepackage{listings}
\lstset{basicstyle=\footnotesize, captionpos=b, breaklines=true, showstringspaces=false, tabsize=2, frame=lines, numbers=left, numberstyle=\tiny, xleftmargin=2em, framexleftmargin=2em}
\makeatletter
\def\l@lstlisting#1#2{\@dottedtocline{1}{0em}{1em}{\hspace{1,5em} Lst. #1}{#2}}
\makeatother

\geometry{a4paper, top=27mm, left=30mm, right=20mm, bottom=35mm, headsep=10mm, footskip=12mm}

\hypersetup{unicode=false, pdftoolbar=true, pdfmenubar=true, pdffitwindow=false, pdfstartview={FitH},
	pdftitle={Entwurf und Implementierung eines Chrome Plugins zur automatischen Anreicherung von Webseiten mit kulturellen Inhalten},
	pdfauthor={Mathias Möller},
	pdfsubject={Bachelorarbeit},
	pdfcreator={\LaTeX\ with package \flqq hyperref\frqq},
	pdfproducer={pdfTeX \the\pdftexversion.\pdftexrevision},
	pdfkeywords={Bachlorarbeit},
	pdfnewwindow=true,
	colorlinks=true,linkcolor=black,citecolor=black,filecolor=magenta,urlcolor=black}
\pdfinfo{/CreationDate (D:20110620133321)}

\begin{document}

\titlespacing{\section}{0pt}{12pt plus 4pt minus 2pt}{-6pt plus 2pt minus 2pt}

% Kopf- und Fusszeile
\renewcommand{\sectionmark}[1]{\markright{#1}}
\renewcommand{\leftmark}{\rightmark}
\pagestyle{fancy}
\lhead{}
\chead{}
\rhead{\thesection\space\contentsname}
\lfoot{Entwurf und Implementierung eines Chrome Plugins zur automatischen\newline Anreicherung von Webseiten mit kulturellen Inhalten}
\cfoot{}
\rfoot{\ \linebreak Seite \thepage}
\renewcommand{\headrulewidth}{0.4pt}
\renewcommand{\footrulewidth}{0.4pt}

% Vorspann
\renewcommand{\thesection}{\Roman{section}}
\renewcommand{\theHsection}{\Roman{section}}
\pagenumbering{Roman}

% ----------------------------------------------------------------------------------------------------------
% Titelseite
% ----------------------------------------------------------------------------------------------------------
\thispagestyle{empty}
\begin{center}
	\includegraphics[scale=0.3]{Bilder/uni_passau.png}\\
	\vspace*{2cm}
	\Large
	\textbf{Fakultät für}\\
	\textbf{Informatik und Mathematik}\\
	\vspace*{2cm}
	\Huge
	\textbf{Bachelorarbeit}\\
	\vspace*{0.5cm}
	\large
	über das Thema\\
	\vspace*{1cm}
	\textbf{Entwurf und Implementierung eines Chrome Plugins zur automatischen Anreicherung von Webseiten mit kulturellen Inhalten}\\
	\vspace*{2cm}
	
	\vfill
	\normalsize
	\newcolumntype{x}[1]{>{\raggedleft\arraybackslash\hspace{0pt}}p{#1}}
	\begin{tabular}{x{6cm}p{7.5cm}}
		\rule{0mm}{5ex}\textbf{Autor:} & Mathias Möller\newline moellerm@fim.uni-passau.de \\ 
		\rule{0mm}{5ex}\textbf{Prüfer:} & Prof. Granitzer \\ 
		\rule{0mm}{5ex}\textbf{Abgabedatum:} & 24.08.2015 \\ 
	\end{tabular} 
\end{center}
\pagebreak

% ----------------------------------------------------------------------------------------------------------
% Abstract
% ----------------------------------------------------------------------------------------------------------
\setcounter{page}{1}
\onehalfspacing
\titlespacing{\section}{0pt}{12pt plus 4pt minus 2pt}{2pt plus 2pt minus 2pt}
\rhead{KURZFASSUNG}
\section{Kurzfassung}
Lorem ipsum dolor sit amet, consetetur sadipscing elitr, sed diam nonumy eirmod tempor invidunt ut labore et dolore magna aliquyam erat, sed diam voluptua. At vero eos et accusam et justo duo dolores et ea rebum. Stet clita kasd gubergren, no sea takimata sanctus est Lorem ipsum dolor sit amet. Lorem ipsum dolor sit amet, consetetur sadipscing elitr, sed diam nonumy eirmod tempor invidunt ut labore et dolore magna aliquyam erat, sed diam voluptua. At vero eos et accusam et justo duo dolores et ea rebum. Stet clita kasd gubergren, no sea takimata sanctus est Lorem ipsum dolor sit amet.

\vspace{-1,2em}
\titlespacing{\section}{0pt}{12pt plus 4pt minus 2pt}{-6pt plus 2pt minus 2pt}
\section*{Abstract}
Das ganze auf Englisch.

\pagebreak

% ----------------------------------------------------------------------------------------------------------
% Verzeichnisse
% ----------------------------------------------------------------------------------------------------------
% TODO Typ vor Nummer
\renewcommand{\cfttabpresnum}{Tab. }
\renewcommand{\cftfigpresnum}{Abb. }
\settowidth{\cfttabnumwidth}{Abb. 10\quad}
\settowidth{\cftfignumwidth}{Abb. 10\quad}

\titlespacing{\section}{0pt}{12pt plus 4pt minus 2pt}{2pt plus 2pt minus 2pt}
\singlespacing
\rhead{INHALTSVERZEICHNIS}
\renewcommand{\contentsname}{II Inhaltsverzeichnis}
\phantomsection
\addcontentsline{toc}{section}{\texorpdfstring{II \hspace{0.35em}Inhaltsverzeichnis}{Inhaltsverzeichnis}}
\addtocounter{section}{1}
\tableofcontents
\pagebreak
\rhead{VERZEICHNISSE}
\listoffigures
\pagebreak
%\listoftables
%\pagebreak

% ----------------------------------------------------------------------------------------------------------
% Abkürzungen
% ----------------------------------------------------------------------------------------------------------
\section{Abkürzungsverzeichnis}
\begin{acronym}[JITIR] % längste Abkürzung steht in eckigen Klammern
	\setlength{\itemsep}{-\parsep} % geringerer Zeilenabstand
	\acro{JITIR}{Just-in-time information retrieval}
\end{acronym}
\newpage

% ----------------------------------------------------------------------------------------------------------
% Inhalt
% ----------------------------------------------------------------------------------------------------------
% Abstände Überschrift
\titlespacing{\section}{0pt}{12pt plus 4pt minus 2pt}{-6pt plus 2pt minus 2pt}
\titlespacing{\subsection}{0pt}{12pt plus 4pt minus 2pt}{-6pt plus 2pt minus 2pt}
\titlespacing{\subsubsection}{0pt}{12pt plus 4pt minus 2pt}{-6pt plus 2pt minus 2pt}

% Kopfzeile
\renewcommand{\sectionmark}[1]{\markright{#1}}
\renewcommand{\subsectionmark}[1]{}
\renewcommand{\subsubsectionmark}[1]{}
\lhead{Kapitel \thesection}
\rhead{\rightmark}

\onehalfspacing
\renewcommand{\thesection}{\arabic{section}}
\renewcommand{\theHsection}{\arabic{section}}
\setcounter{section}{0}
\pagenumbering{arabic}
\setcounter{page}{1}

\section{Einleitung und Motivation}
Suchmaschinen gehören heutzutage zu den meistbesuchten Seiten im Internet. Google\footnote{https://www.google.com/} ist mit 2.000.000.000.000 Suchanfragen pro Jahr (TODO: Quelle) nicht nur die beliebteste Suchmaschine, sondern auch die meist besuchte Seite im Internet. Auch die chinesische Suchmaschine Baidu\footnote{http://www.baidu.com/} ist in den oberen fünf Plätzen vertreten \cite{alexa}. 
\begin{itemize}
	\item Auf XX\% der Webseiten wird über eine Suchmaschine zugegriffen
	\item Menschen holen Informationen über Suchmaschinen, nicht mehr aus Bibliotheken oder Lexika 
	\item Woher Quellen?
\end{itemize}
Um so mehr sie an Bedeutung gewinnen, umso größer ist der Bedarf an bessern Suchmaschinen \cite{lawrence2000context}.

Doch der Funktionsumfang klassischer Suchmaschinen ist begrenzt. Sie können den Kontext des Benutzers in ihren Queries? nicht miteinbeziehen. Weiterhin verlangen sie vom Nutzer, seine gegenwärtige Arbeit einzustellen, die Webseite der Suchmaschine aufzurufen und eine Suchanfrage zu formulieren. So wird die Konzentration auf den eigentlichen Arbeitsschritt gestört. Im Schnitt sind nur 67\% der Suchanfragen an Google erfolgreich (TODO: Quelle). Das heißt in circa einem Drittel der Fälle bleibt der Aufruf der Suchmaschine ohne die erwarteten Erfolge. Zu der verlorenen Zeit addiert sich auch noch die Zeit, die der Nutzer braucht um sich seine Arbeit wieder ins Kurzzeitgedächtnis zu holen.

Um diesen Problemen entgegen zu wirken wurden neue Wege zur Informationsgewinnung entwickelt. Eine Möglichkeit sind so genannte Just-in-time Information Retrieval Agents (JITIR-Agents) \cite{rhodes2000just}. Sie beobachten im Hintergrund den Kontext des Benutzers und versuchen aus den so erhaltenen Informationen eine Suchanfrage an eine Datenbank oder ein Recommender-System zu schicken. Die gewonnenen Informationen werden dann möglichst unaufdringlich dem Benutzer angezeigt. Er kann sich nun entscheiden diese Informationen genauer zu betrachten oder mit seiner Arbeit fortzufahren. Die kognitive Belastung bleibt hierbei sehr gering. JITIR-Agents reduzieren auf diese Weise enorm den Aufwand Informationen zu finden \cite{rhodes2000just}.

Durch ihre Funktionsweise sind sie jedoch nicht so exakt wie klassische Suchmaschinen, da sie nur ``erraten'' können, was den Benutzer gerade interessiert. Wenn ein Benutzer einer genaue Vorstellung von der Suchanfrage oder den Ergebnissen hat, haben klassische Suchmaschinen Vorteile gegenüber den JITIR-Agents \cite{rhodes2000just}.

Ein weiteres Problem des Internets ist es, dass die Standards, die die simple und leicht skalierbare Architektur des Netzwerks ermöglichen, den Einsatz von fortgeschritteneren Hypermedia-Technologien verhindern \cite{bouvin1999unifying}. Eine Möglichkeit, das Internet um Funktionen zu erweitern ist Web Augmentation (WA). Diaz \cite{diaz2012understanding} beschreibt WA als den Versuch, statt eine neue Technologie zu entwickeln, neue Funktionalität auf eine gerenderte Webseite zu setzen:
\begin{quote}
In some sense, WA is to the Web what Augmented Reality is to the physical world: layering relevant content/layout/navigation over the existing Web to customize the user experience. \cite{diaz2012understanding}
\end{quote}
Solche WA-Tools interagieren mit einem Web Server, Http-Proxy oder dem Browser der Nutzers, um so Inhalte oder Navigationselemente direkt in die angezeigte Webseite einzufügen. Auf diese Weise erlauben sie es, die limitierten Möglichkeiten des World Wide Webs zu Gunsten des Nutzers anzureichern (QUelle: Kennth anderson oder Bouvin).

Ein Beispiel für ein solches WA-Tool sowie für ein JITIR System ist das EEXCESS Chrome Plugin\footnote{http://eexcess.eu/results/chrome-extension/}. Es analysiert die aktuell besuchten Seiten des Nutzers und schlägt ihm auf Grund der erlangten Daten am Rand des Browser-Fensters weiterführende Quellen aus der Europeana-Datenbank\footnote{http://www.europeana.eu/portal/} vor.

Ziel dieser Bachelorarbeit ist es, auf Basis des EEXCESS Plugins ein Chrome Plugin zu entwerfen und zu implementieren, welche eine Webseite in einzelne Paragraphen aufteilt und zu jedem dieser Paragraphen kulturelle Inhalte der Europeana-Datenbank vorschlägt. Dabei soll das Design des Plugins helfen, die bei der in Kapitel 5 erläuterten Evaluierung des EEXCESS Plugins aufgetauchten Probleme zu beheben und Schwachstellen zu verbessern.

Im zweiten Teil dieser Arbeit wird auf vergleichbare Technologien eingegangen und ihre Gemeinsamkeiten und Unterschiede zu diesem Projekt. Anschließend wird das Konzept und die Implementierung des Projekts beschrieben und das Ergebnis evaluiert. Der sechste Teil der Ausarbeitung widmet sich Möglichkeiten für zukünftige Projekte und Verbesserungen. Zuletzt wird ein Fazit gezogen und die Ergebnisse der Arbeit analysiert.
\section{Related Work}
Rhodes \cite{rhodes2000just} definiert JITIR-Agents als eine Klasse von Programmen, die dem Benutzer weiterführende Informationen basierend auf seinem lokalen Kontext anzeigen. Dabei beschränkt sich der überwachte Kontext meist auf die virtuelle Umgebung des Benutzers, wie E-Mail, Webseiten und geöffnete Dokumente. Als Kerneigenschaften von JITIR-Agents nennt Rhodes Selbstständigkeit, die Fähigkeit Informationen in einer leicht zugänglichen und gleichzeitig unaufdringlichen Weise darzustellen und das Bewusstsein über den Kontext des Benutzers \cite{rhodes2000thesis}.
Das im Zuge dieser Bachelorarbeit entworfene Programm (im Folgenden als Jarvis bezeichnet) erfordert vom Nutzer die explizite Aufforderung um nach Informationen zu suchen. Gründe für diese Entscheidung werden im dritten Kapitel näher betrachtet. Trotz dieses Widerspruchs zu Rhodes Definition lässt sich Jarvis am besten mit dieser Klasse von Programmen vergleichen. Warum diese Klassifizierung zutrifft wird im folgenden Abschnitt beschrieben.

\subsection{Unterschiede und Gemeinsamkeiten zu JITIR-Agents}
Studien haben gezeigt, dass schon eine kleine Steigerung des Aufwands, der betrieben werden muss um eine Aufgabe zu erfüllen, dazu führen kann, dass man die Aufgabe gar nicht erst ausführt \cite{rhodes2000just}. Laut Miller reicht für die meisten Aufgaben eine Antwortverzögerung von mehr als zwei Sekunden aus, um die Nutzungshäufigkeit des dazugehörigen Programms zu vermindern \cite{miller1968response}. Längere Zeitintervalle erschweren es, den Kontext der gerade ausgeführten Aufgabe und die übergeordneten Aufgaben im Kurzeitgedächtnis zu behalten. Nun wird der Fall betrachtet, dass das Lesen einer Webseite unterbrochen wird, um eine Suche mit einer Suchmaschine durchzuführen. Die eigentliche Aufgabe behält der Nutzer im Kurzzeitgedächtnis. Je länger die Suche dauert und je mehr er sich dazu von seiner eigentlichen Arbeit distanzieren muss, desto schwerer wird es wieder zur Hauptaufgabe zurück zu kehren. Wenn der Exkurs zur Suchmaschine schwerer wiegt ist als die Güte der erwarteten Resultate wird die Suche nicht durchgeführt \cite{rhodes2000just}.

Dieses Problem wird versucht mit JITIR-Agents zu beheben. Durch ihre proaktive Arbeitsweise muss der Nutzer seine Tätigkeit nicht mehr unterbrechen, sondern nur noch entscheiden ob er weiter Informationen sehen möchte oder nicht \cite{rhodes2000just}. Beim Entwurf von Jarvis wurde entschieden, dass das Programm nicht völlig eigenständig nach Informationen sucht, sondern nur die Webseite in seine Paragraphen aufteilt. Der Nutzer kann dann entscheiden, ob er zu einem Paragraphen eine Suche durchführen möchte und diese dann per Klick starten. Wie bei einem JITIR-Agent wird die Suchanfrage automatisch aus den gewonnenen Kontextinformationen generiert. Da die Suche innerhalb weniger Millisekunden ausgelöst werden kann und sich der Nutzer dazu nicht von seiner eigentlichen Aufgabe distanzieren muss, bleibt die kognitive Belastung sehr gering. Allerdings entsteht auch so der Nachteil, den Jarvis mit JITIR-Agents gemein hat: Sie nutzen alle gefunden Informationen für ihre Suchanfragen und können nicht zwischen relevanten und unwichtigen Suchwörtern unterscheiden \cite{rhodes2000margin}. Automatisch gebaute Suchanfragen sind deshalb weniger exakt als Menschen-generierte \cite{rhodes2000just}. Um dem entgegen zu wirken hat der Benutzer von Jarvis im Nachhinein noch die Möglichkeit, die Suche anzupassen und erneut abzuschicken.

Die zweite von Rhodes beschriebene Eigenschaft von JITIR-Agents, die Fähigkeit Informationen in einer leicht zugänglichen und gleichzeitig unaufdringlichen Wiese darzustellen, hat auch beim Entwurf von Jarvis eine bedeutende Rolle gespielt. Für die Darstellung der Ergebnisse muss ein Mittelwert gefunden werden. Die zusätzlichen Elemente sollen den Benutzer nicht unnötig ablenken, allerdings will man die Seite um möglichst reichhaltige Informationen erweitern \cite{rhodes2000margin}. Wie diese Problematik im Falle von Jarvis gelöst wurde wird im Implementierungs-Teil dieser Arbeit beschrieben.

Jarvis analysiert die geöffnete Webseite und teilt diese in Paragraphen ein. Wie JITIR-Agents ist er sich also über den lokalen Kontext des Benutzers bewusst und wertet diesen aus. Die dritte Voraussetzung für JITIR-Agents ist somit erfüllt.

Trotz der eingeschränkten Selbstständigkeit lässt sich Jarvis folglich am besten mit denen von Rhodes beschriebenen Programmen vergleichen. Auf diese Erkenntnis aufbauend wird nachfolgend Jarvis mit existierenden Implementierungen von JITIR-Agents verglichen.

\subsection{Vergleich mit existierenden JITIR-Agents}
Es wurden einige JITIR-Agents in der Vergangenheit implementiert, die zwar aus technischer Hinsicht nicht mehr aktuell sind, die sich jedoch durch ihre genaue Untersuchung und Auswertung gut Vergleichen lassen.

 	\subsubsection{Remembrance Agent}
 	Der Remembrance Agent (RA) ist ein in den UNIX Texteditor EMACS-19 integriertes Programm, welches eine Liste von Dokumenten anzeigt, die für den Nutzer interessant seien könnten \cite{rhodes2000thesis}. Die Informationen für die Vorschläge bezieht der RA aus der gerade geöffneten Datei. Wie Jarvis teilt er das Dokument in Bereiche ein, die er danach analysiert. Allerdings sind die Bereiche in diesem Fall keine Paragraphen sondern verschieden große ``Räume'' in denen er sucht. Von der Position des Cursors ausgehend untersucht er die nächsten zehn, die nächsten 50 und die nächsten 1000 Wörter des Dokuments und sucht zu jedem passende Vorschläge \cite{rhodes1996remembrance}. Die Vorschläge werden von einem Information Retrieval Programm im Hintergrund generiert, welches die Suchanfragen an verschiedene, vom Nutzer konfigurierbare Datenbanken schickt. Genannte Datenbanken sind persönliche Email-Verzeichnisse, persönliche Notizen oder Datenbanken mit wissenschaftlichen Arbeiten.

 	Die gefundenen Ergebnisse werden dem Nutzer dann am unteren Fensterrand dargestellt. Mehrere Zeilen präsentieren die Vorschläge mit der höchsten Relevanz. Je nach Art der Quellen werden verschiedene Informationen angezeigt. So werden zum Beispiel bei wissenschaftlichen Arbeiten der Autor, das Datum und der Titel angezeigt und bei Zeitungsartikeln nur Herausgeber, Überschrift und Datum \cite{rhodes2000just}. Dieses Design ermöglicht eine unaufdringliche Darstellung der Vorschläge, die leicht überblickt werden kann aber den Nutzer nicht von seiner Arbeit ablenkt \cite{rhodes1996remembrance}. Per Klick kann sich der Nutzer die Schlagwörter anzeigen lassen, die für das jeweilige Ergebnis verantwortlich sind \cite{rhodes2000thesis}. Mit einer Tastenkombination und der Zeilennummer des Vorschlags kann er sich das dazugehörende, vollständige Dokument anzeigen lassen. Weiterhin hat der Nutzer die Möglichkeit, eine eigene Suchanfrage einzugeben um so explizit nach Inhalten zu suchen \cite{rhodes1996remembrance}.

 	Mit Hilfe einer Nutzer-Evaluation konnte gezeigt werden, dass der RA nicht nur eine Alternative zu klassischen Suchmaschinen darstellt, sondern in vielen Bereichen sogar besser abscheidet. Durch die Nutzung fanden die Tester mehr relevante Dokumente zum geforderten Themengebiet und die anschließende Umfrage schnitt der RA besser ab als die der Kontrollgruppe zur Verfügung gestellte Suchmaschine \cite{rhodes2000just}.

 	\subsubsection{Watson}
 	Watson ist ein weiteres Beispiel für JITIR-Agents\footnote{Die Entwickler beschreiben das System als ``Information Management Assistant'', die Bedeutung ist jedoch ähnlich.}. Er überwacht die Benutzung von Textverarbeitungs- und Textdarstellungsprogrammen wie Microsoft Word, Microsoft Internet Explorer oder Netscape Navigator \cite{budzik1999watson}. Dazu werden sogenannte ``Application Adapter'' benutzt, die sich mit der jeweiligen Software verbinden um so an den Inhalt der angezeigten oder bearbeiteten Dokumente zu gelangen \cite{budzik2000user}. Die erlangten Informationen werden an Watson weitergegeben, welcher versucht eine passende Quelle aufgrund dieser Daten auszuwählen. Ein weiterer Prozess versucht zu erkennen, ob der Nutzer Bedarf an zusätzlichen Informationen hat. Trifft das zu, wird eine Anfrage an die ausgewählte Quelle/die ausgewählten Quellen geschickt. Die Ergebnisse werden dann gruppiert und in einem separaten Fenster angezeigt \cite{budzik1999watson}.

 	Auch bietet Watson die Möglichkeit, eine Suchanfrage direkt einzugeben. Die Anfrage des Nutzers wird daraufhin mit der automatisch erstellten, kontextabhängigen Anfrage kombiniert um so möglichst relevante Ergebnisse zu liefern \cite{budzik2000user}. Weiterhin werden atomare Bausteine, wie Adressen erkannt. Dem Nutzer wird dann ein Knopf angezeigt, über den er zu einer passenden Darstellung des Bausteins gelangt. Im Falle einer Adresse würde er zum Beispiel zu einer Karte weitergeleitet werden, in der die Adresse angezeigt wird.

 	Evaluationen haben gezeigt, dass die von Watson vorgeschlagenen Dokumente in fünf von zehn Fällen relevant waren. Die Kontrollgruppe, welche die Suche manuell mit Alta Vista\footnote{\url{http://de.wikipedia.org/wiki/AltaVista}} durchführen mussten (Watson benutzte als Quelle auch Alta Vista), kamen im Schnitt auf nur drei relevante Ergebnisse. Schlussendlich erzielte Watson bessere Ergebnisse als die Kontrollgruppe in 15 von 19 Fällen \cite{budzik1999watson}.

 	\subsubsection{EEXCESS Chrome Plugin}
 	Das von EEXCESS\footnote{\url{http://eexcess.eu/}} entwickelte Chrome Plugin\footnote{\url{http://eexcess.eu/results/chrome-extension/}} benutzt JITIR Mechanismen um kulturelle Inhalte der Europeana Datenbank für den Nutzer leichter zugänglich zu machen. Die Software analysiert die aktuelle Webseite oder gegebenenfalls einen selektierten Paragraphen innerhalb dieser Seite und die gewonnenen Daten werden in Form einer Anfrage an Europeana geschickt \cite{schlottererweb}. Die personalisierten Resultate werden dem Benutzer am Rande des Browserfensters in einem separierten Bereich angezeigt.

\subsection{Interactive Query Construction}
Ein System welches Interactive Query Construction verwendet, ist das von Russel et al. entwickelte NITELIGHT \cite{russell2008nitelight}. Es unterstützt Nutzer bei der Erstellung und Anpassung von Suchanfragen auf Basis der SPARQL \cite{sparql} Abfragesprache. Das in Java entwickelte Programm erlaubt es dem Nutzer über eine graphische Oberfläche interaktiv semantische Anfragen zu bauen. Dazu werden ihm graphische Elemente zur Verfügung gestellt, die auf Konstrukten der Abfragesprache basieren \cite{russell2008nitelight}. Die in Jarvis verwendete Interactive Query Construction ist stark vereinfacht und ermöglicht es dem Nutzer lediglich, die automatisch generiert Anfrage zu erweitern oder irrelevante Schlüsselwörter zu entfernen. Jedoch ermöglicht die Nutzung dieser Methodik eine leichte Erweiterbarkeit um Funktionen zur Verbesserung der Ergebnisgüte wie zum Beispiel ``Relevance Feedback''. Darauf wird in Kapitel \ref{sec:futureWork} genauer eingegangen. 

\subsection{Text Retrieval Algorithmen}
 		Term Frequency/Inverse Document Frequency algorithm,
 		Text rank
\section{Konzept}
 \subsection{Warum kein Proactiver JITIR-Agent?}
 	- > API Limitierung und decrease cognitive load 
	- > Benutzer entscheidet ob er weitere Informationen erhalten möchte
	- > Benutzer kann Suchanfrage erst anpassen (Nachteil von Margin Notes (Paper 4)
 \subsection{Anzeige der Ergebnisse}
 \subsection{Erklärung der Such-Anfragen Generierung}
 \subsection{Anpassen der Suchanfrage durch den Nutzer}
 \subsection{Verbesserung der Suchanfrage z.B. durch maschinelles Lernen}

\section{Implementierung}
 \subsection{Verwendung von AngularJS für alle Komponenten des Plugins}
 \subsection{Bau der GUIs}
		- > Darf den Benutzer nicht zu sehr ablenken
		- > Ergebnisse müssen in der Nähe ihrer „Quelle“ angezeigt werden (proximity compatibility pricinple)
		- > Benutzer muss klar zwischen Webseite und Augmentation unterscheiden können
		- > buntes, auffälliges Design
		- > Ramping interface: Mehr Benutzerinteraktionen führen zu mehr angezeigten Informationen (Erklärung der Stages) 
 \subsection{Einbindung der REST-Services}

\section{Evaluierung?}
Aufgaben die Benutzer mit EEXCESS Lösen mussten müssen sie jetzt mit Redesign lösen. Vergleich der Ergebnisse?

\section{Future Work}
 \subsection{Alternative Algorithmen zur Textanalyse}
 \subsection{Implementierung einer automatischen Suchanfragen-Verbesserung durch maschinelles Lernen}
 	\begin{itemize}
 		\item Mehr kontextuelle Informationen miteinbeziehen
 		\item Such-Profil des Nutzers erstellen

 	\end{itemize}
 \subsection{Verbesserung der Ergebnis-Güte}
	\begin{itemize}
		\item durch Query Expansion
 		\item durch Filtern der Ergebnisse (mehr Präzision da Ausbeute bei JITIR nicht so relevant) Clustering
 		\item relevance feedback
	\end{itemize}
\subsection{Anbindung weiterer Quellen}
	Automatische Auswahl der richtigen QUelle laut selecting task relevant sources
 \subsection{Anpassung der Anwendung auf mobile Nutzung}

 \section{Conclusion}
 \begin{itemize}
 		\item Steigerung der Effektivität und Produktivität von wissenschaftlichem Arbeiten
 		\item Starke Effektivitätssteigerung durch Punkte aus Future Work möglich?
 		\item information management assistants embody a vision of a future in which users hardly ever form a query to request information. when an information need arises, a system like watson has already anticipated it and provided relevant information to the user before she is even able to ask for it (budzik, watson)
 		\item survey suggested that users are relatively dissatisfied with the results of their searching experience. THis makes concrete the claim that systems designed to help useres in their information seeking tasks are needed in the world (budzik, watson)

 		JITIR's can be thought of as automatic ``query free'' search engines (rhodes, using physical context)
 	\end{itemize}

% ----------------------------------------------------------------------------------------------------------
% Literatur
% ----------------------------------------------------------------------------------------------------------
\renewcommand\refname{Quellenverzeichnis}
\bibliographystyle{myalpha}
\bibliography{bibo}
\pagebreak

% ----------------------------------------------------------------------------------------------------------
% Anhang
% ----------------------------------------------------------------------------------------------------------
\pagenumbering{Roman}
\setcounter{page}{1}
\lhead{Anhang \thesection}

\begin{appendix}
\section*{Anhang}
\phantomsection
\addcontentsline{toc}{section}{Anhang}
\addtocontents{toc}{\vspace{-0.5em}}


\subsection*{GUI Screenshots}
\label{app:screenshot}
Unterkategorie, die nicht im Inhaltsverzeichnis auftaucht.

\end{appendix}


\newpage
\thispagestyle{empty}
\begin{center}
	\vspace*{5em}
	\huge\textbf{Erklärung}\\
\end{center}
\vspace{2em}
Hiermit versichere ich, dass ich meine Abschlussarbeit selbständig verfasst und keine anderen als die angegebenen Quellen und Hilfsmittel benutzt habe.

\vspace{4em}
\begin{minipage}{\linewidth}
	\begin{tabular}{p{15em}p{15em}}
		Datum: &  .......................................................\\
		& \centering (Unterschrift)\\
	\end{tabular}
\end{minipage}

\end{document}
